\documentclass[../../main.tex]{subfiles}
\begin{document}
\section{单参数子群与指数映射}\label{section:ch20-one-parameter}

设 $G$ 是一个 Lie 群。由于左不变向量场自然地由 $G$ 的群结构定义，人们有理由期望在左不变向量场的流的群法则与 $G$ 中的群乘法之间找到某种关系。我们就从探讨这一关系开始。

\subsection{单参数子群}

$G$ 的一个\emph{单参数子群}定义为一个 Lie 群同态 $\gamma\colon\mathbb{R}\to G$，其中 $\mathbb{R}$ 视为加法下的 Lie 群。按此定义，单参数子群并不是 $G$ 的 Lie 子群，而是映到 $G$ 的一个同态。（然而，当赋予适当的光滑流形结构时，单参数子群的像是一个 Lie 子群；见习题 20-1。）

\begin{theorem}{[单参数子群的刻画]}\label{theorem:ch20-单参数子群刻画}
设 $G$ 是一个 Lie 群。$G$ 的单参数子群恰为从单位元出发的左不变向量场的极大积分曲线。
\end{theorem}

\begin{proof}
首先设 $\gamma$ 是某个左不变向量场 $X\in\operatorname{Lie}(G)$ 从单位元出发的极大积分曲线。由于左不变向量场是完全的（定理\cref{theorem:ch09-18}），$\gamma$ 在整个 $\mathbb{R}$ 上有定义。左不变性意味着对每个 $g\in G$，$X$ 与它自身是 $L_g$-相关的，故命题\cref{proposition:ch09-06} 蕴含 $L_g$ 把 $X$ 的积分曲线映为 $X$ 的积分曲线。取 $g=\gamma(s)$（$s\in\mathbb{R}$）应用此结论，可知曲线 $t\mapsto L_{\gamma(s)}\bigl(\gamma(t)\bigr)$ 是从 $\gamma(s)$ 出发的积分曲线。但平移引理（引理\cref{lemma:ch09-04}）蕴含 $t\mapsto\gamma(s+t)$ 也是具有相同初值的积分曲线，故二者相等：
\begin{align*}
\gamma(s)\gamma(t)=\gamma(s+t).
\end{align*}
这恰说明 $\gamma\colon\mathbb{R}\to G$ 是一个单参数子群。

反之，设 $\gamma\colon\mathbb{R}\to G$ 是一个单参数子群，令 $X=\gamma_*(\mathrm{d}/\mathrm{d}t)\in\operatorname{Lie}(G)$，其中把 $\mathrm{d}/\mathrm{d}t$ 视为 $\mathbb{R}$ 上的左不变向量场。由于 $\gamma(0)=e$，只需证明 $\gamma$ 是 $X$ 的积分曲线。回忆 $\gamma_*(\mathrm{d}/\mathrm{d}t)$ 定义为 $G$ 上与 $\mathrm{d}/\mathrm{d}t$ 是 $\gamma$-相关的唯一左不变向量场（定理\cref{theorem:ch08-44}）。因此对任意 $t_0\in\mathbb{R}$，
\begin{align*}
\gamma'(t_0)=\mathrm{d}\gamma_{t_0}\left(\frac{\mathrm{d}}{\mathrm{d}t}\bigg|_{t_0}\right)=X_{\gamma(t_0)},
\end{align*}
故 $\gamma$ 是 $X$ 的积分曲线。
\end{proof}

给定 $X\in\operatorname{Lie}(G)$，如此由 $X$ 确定的单参数子群称为由 $X$ \emph{生成的}单参数子群。由于左不变向量场由其在单位元处的值唯一确定，由此可知每个单参数子群都是由唯一的左不变向量场生成的。

$\operatorname{GL}(n,\mathbb{R})$ 的单参数子群不难显式计算。

\begin{proposition}\label{proposition:ch20-矩阵指数}
对任意 $A\in\mathfrak{gl}(n,\mathbb{R})$，令
\begin{align}
e^A=\sum_{k=0}^\infty \frac1{k!}A^k=I_n+A+\tfrac12A^2+\cdots.
\end{align}
该级数收敛到一个可逆矩阵 $e^A\in\operatorname{GL}(n,\mathbb{R})$，并且由 $A\in\mathfrak{gl}(n,\mathbb{R})$ 生成的 $\operatorname{GL}(n,\mathbb{R})$ 的单参数子群为 $\gamma(t)=e^{tA}$。
\end{proposition}

\begin{proof}
首先验证收敛性。由习题 B.48，矩阵乘法满足 $|AB|\leqslant|A||B|$，其中范数是 $\mathfrak{gl}(n,\mathbb{R})$ 上的 Frobenius 范数。由归纳法得 $|A^k|\leqslant|A|^k$。Weierstrass $M$-判别法表明上式在 $\mathfrak{gl}(n,\mathbb{R})$ 的任意有界子集上一致收敛，因为它与级数 $\sum_k(1/k!)c^k=e^c$ 比较。

固定 $A\in\mathfrak{gl}(n,\mathbb{R})$。在把 $\mathfrak{gl}(n,\mathbb{R})$ 与 $\operatorname{Lie}\bigl(\operatorname{GL}(n,\mathbb{R})\bigr)$ 等同之下，矩阵 $A$ 对应于由\eqref{align:ch08-15} 给出的左不变向量场 $A^L$。于是由 $A$ 生成的单参数子群是 $A^L$ 在 $\operatorname{GL}(n,\mathbb{R})$ 上的积分曲线，从而满足 ODE 初值问题
\begin{align*}
\gamma'(t)=A^L\big|_{\gamma(t)},\qquad \gamma(0)=I_n.
\end{align*}
利用\eqref{align:ch08-15}，$\gamma$ 为积分曲线的条件可改写为 $\dot\gamma^i_k(t)=\gamma^i_j(t)A^j_k$，或写成矩阵形式 $\gamma'(t)=\gamma(t)A$。我们将证明 $\gamma(t)=e^{tA}$ 满足此方程。由于 $\gamma(0)=I_n$，这蕴含 $\gamma$ 是从单位元出发的 $A^L$ 的唯一积分曲线，因而是所求的单参数子群。

为看出 $\gamma$ 可微，注意对级数逐项形式求导得到
\begin{align*}
\gamma'(t)=\sum_{k=1}^\infty \frac1{(k-1)!}t^{k-1}A^k=\left(\sum_{j=0}^\infty\frac1{j!}t^jA^j\right)A=e^{tA}A=\gamma(t)A.
\end{align*}
由此 $\gamma$ 是常值曲线 $\gamma(t)\gamma(-t)=I_n$，即 $\gamma(t)\gamma(-t)=I_n$。把 $t$ 换成 $-t$，得 $\gamma(t)\gamma(-t)=I_n$，这表明 $\gamma(t)$ 可逆且 $\gamma(t)^{-1}=\gamma(-t)$。
\end{proof}

接下来我们想计算 $\operatorname{GL}(n,\mathbb{R})$ 的单参数子群，例如 $O(n)$ 的。为此需要下列结果。

\begin{proposition}\label{proposition:ch20-子群单参数}
设 $G$ 是一个 Lie 群，$H\subseteq G$ 是一个 Lie 子群。$H$ 的单参数子群恰为 $G$ 中那些初速度位于 $T_eH$ 的单参数子群。
\end{proposition}

\begin{proof}
设 $\gamma\colon\mathbb{R}\to H$ 是一个单参数子群。则复合映射 $\mathbb{R}\to H\hookrightarrow G$ 是 Lie 群同态，因而是 $G$ 的一个单参数子群，且显然满足 $\gamma'(0)\in T_eH$。

反之，设 $\gamma\colon\mathbb{R}\to G$ 是一个初速度位于 $T_eH$ 的单参数子群。令 $\widetilde\gamma\colon\mathbb{R}\to H$ 为 $H$ 中与 $\gamma$ 有相同初速度 $\widetilde\gamma'(0)=\gamma'(0)\in T_eH\subseteq T_eG$ 的单参数子群。如前一段那样，通过与包含映射复合，也可把 $\widetilde\gamma$ 视为 $G$ 的单参数子群。由于 $\gamma$ 与 $\widetilde\gamma$ 都是 $G$ 中初速度相同的单参数子群，它们必相等。
\end{proof}

\begin{example}\label{example:ch20-正交群单参数}
若 $H$ 是 $\operatorname{GL}(n,\mathbb{R})$ 的 Lie 子群，上述命题表明 $H$ 的单参数子群恰为形如 $\gamma(t)=e^{tA}$（$A\in\mathfrak h$）的映射，其中 $\mathfrak h\subseteq\mathfrak{gl}(n,\mathbb{R})$ 是按定理\cref{theorem:ch08-46} 与 $\operatorname{Lie}(H)$ 对应的子代数。例如取 $H=O(n)$，这表明 $O(n)$ 的单参数子群是对任意反对称矩阵 $A$ 形如 $\gamma(t)=e^{tA}$ 的映射。特别地，这表明任意反对称矩阵的指数都是正交矩阵。
\end{example}

\subsection{指数映射}

上一节我们看到矩阵指数把 $\mathfrak{gl}(n,\mathbb{R})$ 映到 $\operatorname{GL}(n,\mathbb{R})$，并把每条过原点的直线映为一个单参数子群。这可强有力地推广到任意 Lie 群。

给定 Lie 群 $G$ 及其 Lie 代数 $\mathfrak g$，我们定义映射 $\exp\colon\mathfrak g\to G$，称为 $G$ 的\emph{指数映射}，如下：对任意 $X\in\mathfrak g$，令
\begin{align*}
\exp X=\gamma(1),
\end{align*}
其中 $\gamma$ 是由 $X$ 生成的单参数子群，等价地说是从单位元出发的 $X$ 的积分曲线（见图\cref{figure:ch20-exponential}）。下列命题表明，与矩阵指数一样，该映射把过 $X$ 的直线映为由 $X$ 生成的单参数子群。

\begin{figure}[ht]
\centering
\includegraphics[width=0.8\textwidth]{fig-20-1.png}
\caption{指数映射}\label{figure:ch20-exponential}
\end{figure}

\begin{proposition}\label{proposition:ch20-exp直线}
设 $G$ 是一个 Lie 群。对任意 $X\in\operatorname{Lie}(G)$，$\gamma(s)=\exp sX$ 是由 $X$ 生成的 $G$ 的单参数子群。
\end{proposition}

\begin{proof}
设 $\gamma\colon\mathbb{R}\to G$ 是由 $X$ 生成的单参数子群，即从 $e$ 出发的 $X$ 的积分曲线。对任意固定的 $s\in\mathbb{R}$，由重标度引理（引理\cref{lemma:ch09-03}），$\widetilde\gamma(t)=\gamma(st)$ 是从 $e$ 出发的 $sX$ 的积分曲线，故 $\exp sX=\widetilde\gamma(1)=\gamma(s)$。
\end{proof}

下面是两个简单而重要的例子。

\begin{example}\label{example:ch20-gl指数}
上一节的结果表明 $\operatorname{GL}(n,\mathbb{R})$（或其任意 Lie 子群）的指数映射由 $\exp A=e^A$ 给出。显然这正是"指数映射"这一名称的由来。
\end{example}

\begin{example}\label{example:ch20-向量空间指数}
若 $V$ 是有限维实向量空间，选取 $V$ 的一组基给出同构 $\operatorname{GL}(V)\cong\operatorname{GL}(n,\mathbb{R})$ 与 $\mathfrak{gl}(V)\cong\mathfrak{gl}(n,\mathbb{R})$。对 $\operatorname{GL}(n,\mathbb{R})$ 情形的分析表明 $\operatorname{GL}(V)$ 的指数映射可写成
\begin{align}
\exp A=\sum_{k=0}^\infty\frac1{k!}A^k,
\end{align}
其中把 $A\in\mathfrak{gl}(V)$ 视为从 $V$ 到自身的线性映射，$A^k=A\circ\cdots\circ A$ 是 $A$ 与自身的 $k$ 重复合。
\end{example}

\begin{proposition}{[指数映射的性质]}\label{proposition:ch20-指数性质}
设 $G$ 是一个 Lie 群，$\mathfrak g$ 是它的 Lie 代数。
\begin{enumerate}[(a)]
\item 指数映射是从 $\mathfrak g$ 到 $G$ 的光滑映射。
\item 对任意 $X\in\mathfrak g$ 与 $s,t\in\mathbb{R}$，$\exp(s+t)X=(\exp sX)(\exp tX)$。
\item 对任意 $X\in\mathfrak g$，$(\exp X)^{-1}=\exp(-X)$。
\item 对任意 $X\in\mathfrak g$ 与 $n\in\mathbb{Z}$，$(\exp X)^n=\exp(nX)$。
\item 微分 $(\mathrm{d}\exp)_0\colon T_0\mathfrak g\to T_eG$ 在把 $T_0\mathfrak g$ 与 $T_eG$ 都典范等同于 $\mathfrak g$ 之下是恒等映射。
\item 指数映射限制为从 $\mathfrak g$ 中 $0$ 的某个邻域到 $G$ 中 $e$ 的某个邻域的微分同胚。
\item 若 $H$ 是另一个 Lie 群，$\mathfrak h$ 是它的 Lie 代数，$\Phi\colon G\to H$ 是 Lie 群同态，则图表
\begin{align*}
\begin{CD}
\mathfrak g @>{\Phi_*}>> \mathfrak h\\
@V{\exp}VV @VV{\exp}V\\
G @>{\Phi}>> H
\end{CD}
\end{align*}
交换。
\item 左不变向量场 $X$ 的流 $\theta^{(X)}$ 由 $\theta^{(X)}_t=R_{\exp tX}$（即右乘 $\exp tX$）给出。
\end{enumerate}
\end{proposition}

\begin{proof}
在本证明中，对任意 $X\in\mathfrak g$，令 $\theta^{(X)}$ 表示 $X$ 的流。为证 (a)，需表明 $\gamma_e(X)=\theta^{(X)}(1)$ 光滑依赖于 $X$，这归结为表明流随向量场光滑地变化。这是流的基本定理未涵盖的情形，但可用下面的简单技巧化归到该定理。在乘积流形 $G\times\mathfrak g$ 上定义向量场
\begin{align*}
\Gamma(g,X)=(X_g,0)\in T_gG\oplus T_X\mathfrak g\cong T_{(g,X)}(G\times\mathfrak g).
\end{align*}
为看出 $\Gamma$ 是光滑向量场，取 $\mathfrak g$ 的一组基 $(X_1,\dots,X_k)$，令 $(\xi^i)$ 为 $\mathfrak g$ 上由 $\sum\xi^iX_i$ 给出的对应整体坐标，$(w^i)$ 为 $G$ 的任意光滑局部坐标。若 $f\in C^\infty(G\times\mathfrak g)$ 任意，则局部可写成 $f(w^i,\xi^i)=\sum_j\xi^jX_jf(w^i,\xi^i)$，其中每个 $X_j$ 只在 $w^i$ 方向上对 $f$ 求导。由于这光滑依赖于 $(\xi^i)$，$\Gamma$ 光滑。

$\Gamma$ 的积分曲线在 $G$ 分量上的投影恰为 $X$ 的流；又因其 $\mathfrak g$ 分量为常数，故在时刻 $1$，初始值为 $(e,X)$ 的积分曲线其 $G$ 分量等于从 $e$ 出发的 $X$ 的积分曲线在时刻 $1$ 的值，即 $\exp X$。由光滑依赖参数，$\exp X$ 光滑依赖于 $X$。这就证明了 (a)。

(b) 由命题\cref{proposition:ch20-exp直线}，$t\mapsto\exp tX$ 是单参数子群，故 $\exp(s+t)X=(\exp sX)(\exp tX)$。

(c) 由 (b)，$(\exp X)(\exp(-X))=\exp 0=e$，故 $(\exp X)^{-1}=\exp(-X)$。

(d) 由 (b) 对正整数 $n$ 归纳，并由 (c) 过渡到负整数。

(e) 把 $\exp$ 视为 $t\mapsto\exp tX$ 在 $t=0$ 处关于 $t$ 的导数，并注意这是从 $e$ 出发的 $X$ 的积分曲线，其初速度恰为 $X_e$；在典范等同下 $(\mathrm{d}\exp)_0(X)=X$。

(f) 由 (e) 与反函数定理。

(g) 设 $\gamma(t)=\exp tX$ 是由 $X$ 生成的单参数子群。左端按定义是由 $\Phi_*X$ 生成的单参数子群，即在 $H$ 中从 $e$ 出发的 $\Phi_*X$ 的积分曲线。于是若令 $\widetilde\gamma(t)=\Phi(\exp tX)$，只需证明 $\widetilde\gamma\colon\mathbb{R}\to H$ 是满足 $\widetilde\gamma'(0)=(\Phi_*X)_e$ 的 Lie 群同态。它是 Lie 群同态，因为它是同态 $\Phi$ 与 $t\mapsto\exp tX$ 的复合。初速度计算如下：
\begin{align*}
\widetilde\gamma'(0)=\frac{\mathrm{d}}{\mathrm{d}t}\bigg|_{t=0}\Phi(\exp tX)=\mathrm{d}\Phi_e\left(\frac{\mathrm{d}}{\mathrm{d}t}\bigg|_{t=0}\exp tX\right)=\mathrm{d}\Phi_e(X_e)=(\Phi_*X)_e.
\end{align*}

(h) 为表明 $\theta^{(X)}_t=R_{\exp tX}$，利用对任意 $g\in G$，左乘映射 $L_g$ 把 $X$ 的积分曲线映为 $X$ 的积分曲线。于是映射 $t\mapsto L_g(\exp tX)$ 是从 $g$ 出发的积分曲线，即等于 $\theta^{(X)}_t(g)$。由此
\begin{align*}
R_{\exp tX}(g)=g\exp tX=L_g(\exp tX)=\theta^{(X)}_t(g)=\theta^{(X)}_t(g).
\end{align*}
\end{proof}

需要注意，命题\cref{proposition:ch20-指数性质}(b) 并不蕴含对 Lie 代数中任意 $X,Y$ 有 $\exp(X+Y)=(\exp X)(\exp Y)$。事实上，对连通群而言，这只在群交换时才成立（见习题 20-8）。

指数映射给出子群的 Lie 代数的另一种刻画。我们在本章后面研究正规子群时会用到它。

\begin{proposition}\label{proposition:ch20-子群Lie代数}
设 $G$ 是一个 Lie 群，$H\subseteq G$ 是一个 Lie 子群。按通常方式把 $\operatorname{Lie}(H)$ 视为 $\operatorname{Lie}(G)$ 的子代数，则 $H$ 的指数映射是 $G$ 的指数映射在 $\operatorname{Lie}(H)$ 上的限制，并且
\begin{align*}
\operatorname{Lie}(H)=\{X\in\operatorname{Lie}(G)\colon \exp tX\in H\text{ 对所有 }t\in\mathbb{R}\}.
\end{align*}
\end{proposition}

\begin{proof}
$H$ 的指数映射是 $G$ 的指数映射的限制，这是命题\cref{proposition:ch20-子群单参数} 的直接推论。为证第二个结论，按我们把 $\operatorname{Lie}(H)$ 等同于 $\operatorname{Lie}(G)$ 的子代数的方式，需对每个 $X\in\operatorname{Lie}(G)$ 建立等价关系：$\exp tX\in H$ 对所有 $t\in\mathbb{R}$ 当且仅当 $X_e\in T_eH$。

首先设 $\exp tX\in H$ 对所有 $t$ 成立。由于 $H$ 由定理\cref{theorem:ch19-子群弱嵌入} 在 $G$ 中弱嵌入，曲线 $t\mapsto\exp tX$ 作为到 $H$ 的映射是光滑的，故 $X_e=\gamma'(0)\in T_eH$。反之，若 $X_e\in T_eH$，则命题\cref{proposition:ch20-子群单参数} 蕴含 $\exp tX\in H$ 对所有 $t$ 成立。
\end{proof}
\end{document}
